\chapter{Countable tests for distributional currents}\label{sec:current-test-density}
Slicing produces order-one currents before the finite-mass estimate is known. Countable tests must therefore approximate first derivatives as well as coefficients. The following functional-analytic step expands the density argument used in \cite{Sim18}, Chapter 6, p.~183.

\begin{lemma}[Countable tests in the first-derivative topology]\label{lem:countable-current-c1-tests}
For every open $W\subset\mathbb R^N$ and form degree $q\geq0$, there is a countable set of smooth compactly supported $q$-forms in $W$ such that every smooth test is approximated in the $C^1$ norm by a sequence from this set, with all sequence supports in one compact subset of $W$. There is also such a countable set within the comass unit ball which is dense among tests with comass at most one. See the countable-density step in \cite{Sim18}, Lemma 4.5.
\end{lemma}
\begin{proof}
Fix a nonnegative smooth compactly supported scalar convolution kernel. Extend each test coefficient by zero to $\mathbb R^N$ and convolve at rational positive scales tending to zero. A smooth compactly supported function and its first derivatives converge uniformly under this operation; for sufficiently small scales all supports remain in a fixed compact subset of $W$. For a fixed scale, approximate the convolution integral and all its first derivatives by finite Riemann sums with rational centers. Uniform continuity of the kernel and its first derivatives, together with compact integration support, justifies both approximations simultaneously. Replace the finitely many real coefficients of the sums by rational numbers, changing the $C^1$ error by an arbitrarily small amount. Finite rational combinations of rational-scale kernels centered at rational points form a countable collection. Retain only combinations supported in $W$, and apply this to each component of the form.

For the unit-comass assertion, also allow division by positive rational constants. Given an approximant $\eta_j$ to a unit-comass test, uniform convergence implies $\sup_x\|\eta_j(x)\|_*\leq1+o(1)$. Choose a rational divisor $r_j\geq\max\{1,\sup_x\|\eta_j(x)\|_*\}$ with $r_j\to1$. Then $\eta_j/r_j$ has comass at most one and still converges in $C^1$ on the common compact support. This preserves countability. Add zero and negatives if needed.
\end{proof}

\begin{proposition}[Mass as a countable supremum for order-one currents]\label{prop:order-one-countable-mass}
Let $B$ be a $q$-current on $W$ satisfying, for each compact $K\subset W$, an estimate $|B(\omega)|\leq C_K\|\omega\|_{C^1}$ for tests supported in $K$. For the countable unit-comass family $\{\omega_j\}$ of Lemma \ref{lem:countable-current-c1-tests}, chosen to include zero and negatives,
\[
 \mathbf M_W(B)=\sup_j B(\omega_j),
\]
allowing infinite mass. If $t\mapsto B_t(\omega)$ is measurable for every test and each $B_t$ has this order-one property, then $t\mapsto\mathbf M_W(B_t)$ is measurable. This formalizes the testing step in the slicing estimates of \cite{Sim18}, Section 4.
\end{proposition}
\begin{proof}
Every chosen test is admissible in the definition of mass, giving one inequality. Conversely approximate any unit-comass test by the family in $C^1$ with common compact support. The assumed estimate implies convergence of its evaluations, so its value is bounded by the countable supremum. Taking the supremum over all tests proves the reverse inequality. The supremum of countably many measurable functions is measurable, which proves the last assertion. No uniformity in the constants $C_K$ as $t$ varies is required.
\end{proof}
